\section{BFS}

\subsection{Домашнее задание}
\begin{enumerate}
  \item
    Дан взвешенный орграф с положительными весами и вершина $s$.
    Нужно для каждой вершины $v$ найти число кратчайших путей из $s$ в $v$ по модулю $p$. $\O(E \log V)$.

  \item
    В орграфе есть бесплатные и $K$ типов платных ребер. Передвижение
    по любому ребру стоит $0$, но, чтобы двигаться по платному ребру
    типа $x$, необходимо иметь пропуск того же типа. В любой момент
    времени можно иметь не более одного пропуска, однако в любой
    вершине есть возможность купить пропуск любого типа за $A$ и
    продать пропуск любого типа за $B$ ($0 < B < A$, цены одинаковы
    во всех вершинах). Найдите самый дешевый способ добраться из вершины
    $s$ до вершины $t$ за $\O(K (V + E))$.

  \item
    Дан ориентированный граф с выделенными вершинами $s$ и $t$. Перемещение по любому ребру занимает 1 единицу времени.
    В начальный момент времени мы стоим в вершине $s$, и во все ребра с постоянной скоростью $L$ начинает поступать вода.
    У каждого ребра $e$ есть емкость $C_e$ --- минимальное количество воды в этом ребре, при котором по нему больше
    нельзя двигаться (мы должны закончить движение по ребру не позднее момента достижения его емкости). Найдите
    максимальное $L$, при котором все еще существует возможность добраться от $s$ до $t$ за $\O((V + E)\log{C_{\max}})$.

  \item
    Постройте матрицу $n \times n$, состоящую из клеток-стенок и пустых клеток,
    на которой при запуске \texttt{BFS} из какой-то клетки максимальный
    размер очереди будет $\omega(n)$. Ходить можно между клетками, смежными по стороне.

\subsection*{Дополнительные задачи}

  \item
    Попробуем научить алгоритм Дейкстры работать на графах с отрицательными весами: вынимаем из кучи вершину $v$
    с минимальным расстоянием и релаксируем все ребра из $v$. При этом вершины, до которых улучшилось расстояние,
    добавляются в кучу (или делается \texttt{DecreaseKey}, если вершина уже там).
    \begin{enumerate}
      \item Покажите, что этот алгоритм корректно вычисляет кратчайшие пути, если нет циклов отрицательного веса.
      \item Придумайте пример без циклов отрицательного веса, на котором время работы экспоненциально.
    \end{enumerate}

  \item
    Модифицируем условия задачи про бочки: есть $n$ бочек, $i$-я бочка имеет
    объем $c_i$ ($c_i > 0$) и в начальный момент содержит $x_i$ воды
    ($0 \le x_i \le c_i$). Операция переливания определена как в задаче из практики.
    Придумайте такие $n$ и вещественные $c_i$, $x_i$, чтобы соответствующий граф
    состояний, достижимых с помощью переливаний из начального, получился бесконечным.

  \item
    Дан набор степеней $d_1, d_2, \cdots, d_n$. Требуется
    построить граф $G = \langle V, E \rangle$, чтобы степень вершины
    $v_i$ была равна $d_i$, или сказать, что такого не существует.
    Граф не должен содержать петли и кратные ребра. Решить за $\O(V + E)$.

\end{enumerate}

\pagebreak

\section{Решения}

\begin{enumerate}

    \item
        Для решения задачи воспользуемся алгоритмом Дейкстры с использованием очереди с приоритетом.

        Пусть $d[v]$ --- длина кратчайшего пути из $s$ в $v$, а $c[v]$ --- количество кратчайших путей из $s$ в $v$.
        Тогда $d[s] = 0$, $c[s] = 1$, а для всех остальных вершин $v$: $d[v] = \infty$, $c[v] = 0$.

        Далее запустим алгоритм Дейкстры и на каждом шаге будем извлекать минимальную вершину $v$ из очереди.
        Для каждого ребра $(v, u)$, ведущего в вершину $u$, будем релаксировать его:

        \begin{itemize}
            \item Если $d[v] + w(v, u) < d[u]$, то обновляем $d[u]$:
            \[
            d[u] = d[v] + w(v, u), \quad c[u] = c[v].
            \]
            \item Если $d[v] + w(v, u) = d[u]$, то добавляем количество путей:
            \[
            c[u] = c[u] + c[v].
            \]
            \item Если $d[v] + w(v, u) > d[u]$, то ребро $(v, u)$ не улучшает путь и пропускаем его.
        \end{itemize}

        Таким образом, после завершения работы алгоритма $c[v]$ будет содержать количество кратчайших путей из $s$ в $v$.
        Сложность алгоритма --- $\O(E \log V)$.

    \item

        Для решения построим граф $G$ в котором каждое ребро начального графа будет представлять из себя $K + 1$
        ребро в новом графе с соответствующим состоянием пропусков.
        Где $K$ --- количество типов платных ребер, которые обязывают иметь пропуск соответствующего типа и $1$ --- бесплатное ребро,
        доступное в любой момент времени.

        Для этого используем поиск в глубину с сохранением состояния пропуска, номера вершины и стоимости пути до нее.
        Изначально запускаем поиск из вершины $s$ с состоянием пропуска $0$ и стоимостью $0$.
        На каждой итерации выполняем следующие действия:
        \begin{enumerate}

        \item Бесплатное ребро. Переходим в вершину с тем же состоянием пропуска и стоимостью пути.
        \item Платное ребро. Если у нас есть пропуск соответствующего типа, то переходим в вершину с новым состоянием пропуска и стоимостью пути.
        \item Покупка пропуска. Если у нас нет пропуска, то покупаем его и откладываем событие в очередь из той же вершины с новым состоянием пропуска и стоимостью пути увеличенной на $A$.
        \item Продажа пропуска. Если у нас есть какой-то другой пропуск, то продаем его и откладываем событие в очередь из той же вершины с новым состоянием пропуска и стоимостью пути уменьшенной на $B$.

        \end{enumerate}

        Обрабатывая, каждую вершину мы проверяем, что она минимизует стоимость пути до нее.
        После завершения поиска в глубину, находим минимальную из стоимостей путей до вершины $t$.
        Так как мы увеличили количество вершин и ребер в $K$ раз, то сложность алгоритма --- $\O(K(V + E))$.

    \item

        Для решения задачи используем бинарный поиск по ответу. Нам необходимо найти максимальную скорость поступления воды
        $L$, при которой существует путь из вершины $s$ в вершину $t$. Заметим, что с увеличением $L$ время заполнения ребер
        уменьшается, и при слишком большом $L$ путь из $s$ в $t$ становится невозможным, так как ребра блокируются
        раньше, чем мы успеем пройти через них.

        Введем следующий алгоритм:

        \begin{enumerate}

        \item Границы поиска и обновление скорости.

        Задаем начальные границы поиска: $L_{min} = 0, L_{max} = C_{max}$, где $C_{max}$ --- максимальная емкость среди всех ребер графа.
        На каждой итерации бинарного поиска рассчитываем скорость как $L_{mid} = \frac{L_{min} + L_{max}}{2}$.
        Если путь из $s$ в $t$ существует при скорости $L_{mid}$, то увеличиваем $L_{min}$, иначе уменьшаем $L_{max}$.

        \item Проверка существования пути при скорости $L$.

        Для проверки существования пути из $s$ в $t$ при скорости $L$ запускаем поиск в ширину из вершины $s$.
        При этом для каждой вершины $v$ будем хранить время $t_v$, когда вода достигнет вершины $v$.
        В начале $t_s = 0$, а для всех остальных вершин $v$ $t_v = \infty$.

        Рассматривая ребро $e = (u, v)$
        с емкостью $C_e$, определяем момент времени, когда ребро заблокируется:
        $t_e = \lceil \frac{C_e}{L} \rceil$.
        Если текущий момент времени $t_u + 1 \leq t_e$, то мы можем пройти через ребро $e$,
        и время достижения вершины $v$ обновляется как $t_v = t_u + 1$.

        Если после завершения BFS время достижения вершины $t_t$ осталось равным $\infty$, то путь из $s$ в $t$ не существует.
        Иначе путь существует при скорости $L$.

        \item Условие остановки.

        Бинарный поиск делаем по какой-то дельте, например, $10^{-6}$ и продолжаем до тех пор
        пока $L_{max} - L_{min} > \delta$.

        \end{enumerate}

        Основная идея заключается в том, что скорость воды $L$
        прямо пропорционально влияет на время заполнения ребер. Чем больше
        $L$, тем быстрее вода заполняет ребра, и, следовательно, тем быстрее они блокируются.

        Сложность алгоритма --- $\O((V + E)\log{C_{\max}})$.


    \item \dots
        
\end{enumerate}
